% When player-level accuracy does not reach the squad.
% Compile from this folder with tectonic (XeTeX):
%   tectonic -Z shell-escape=false squadopt_walk_forward_paper.tex
% Figures are regenerated by scripts.paper_figures and converted to PDF under
% artifacts/paper/figures-pdf (git-ignored); see docs/paper/README.md.
\documentclass[11pt,a4paper]{article}

\usepackage{fontspec}
\usepackage{amsmath,amssymb}
\usepackage{booktabs}
\usepackage{graphicx}
\usepackage{caption}
\usepackage{subcaption}
\usepackage{tabularx}
\usepackage{array}
\usepackage{enumitem}
\usepackage{xcolor}
\usepackage[a4paper,margin=27mm]{geometry}
\usepackage{microtype}
\usepackage[hidelinks]{hyperref}

\graphicspath{{../../artifacts/paper/figures-pdf/}}
\captionsetup{font=small,labelfont=bf}
\setlist{nosep}
\newcolumntype{R}{>{\raggedleft\arraybackslash}X}
\newcolumntype{L}{>{\raggedright\arraybackslash}X}
\newcommand{\gw}{\textsc{gw}}
\newcommand{\ci}[2]{$[#1,\,#2]$}

\title{When Player-Level Accuracy Does Not Reach the Squad:\\
A Pre-Registered Walk-Forward Study of Component Forecasts,\\
Decision Value and Scenario Calibration in Fantasy Premier League}

\author{Ertuğrul Soydal\thanks{Both authors contributed equally to this work. Corresponding author: Ertuğrul Soydal, Sabancı University, \texttt{ertugrul.soydal@sabanciuniv.edu}.}\\
Sabancı University, Istanbul
\and İbrahim Ersan Özdemir\footnotemark[1]}
\date{Draft of 8 November 2026}

\begin{document}
\maketitle

\begin{abstract}
Fantasy Premier League (FPL) is a weekly sequential decision problem: a manager holds fifteen players under a budget, a three-per-club cap and positional quotas, names a starting eleven and a captain, and may change the squad between deadlines at a price. A forecasting model enters this problem only through the expected points it hands to a solver, so a better forecast is worth exactly as much as the better squads it produces. We report a pre-registered walk-forward study, on five public seasons of player-gameweek data (2021-22 to 2025-26, 184 folds, 130{,}095 player-gameweek rows), of how far gains in player-level accuracy carry into realized squad points when the squad is chosen by an exact integer programme under the official rules. The forecasting model decomposes a player's gameweek into an appearance probability, conditional minutes and conditional points; the served model fits these with a standardized logistic regression and two ridge regressions on nine as-of features. Five declared candidates replaced its estimators or its features: a histogram gradient-boosted pair, a wider feature contract, boosting of the points component alone, regime-gated experts and a gated two-expert mixture. Two of the five were preferred under the pre-declared points-error rule, with 2025-26 points MAE falling from 1.0058 to 0.9872 and 0.9679; yet on the 147 frozen decision folds, where every arm's squads were solved by the same CP-SAT model and scored under the official automatic-substitution and captain rules, no candidate separated from the linear control: paired means ranged from $-0.95$ to $+0.54$ points a gameweek and every 90\% season-aware block-bootstrap interval covered zero. The two readings disagreed in sign for four of the five candidates. The mixture, no better as a mean, gave a sharper and better-calibrated predictive distribution on every fold scored (points log score $+0.38$ to $+0.51$ a fold), so distribution quality is a separate axis from decision value. For squad-level risk we measured three scenario samplers on two frozen calibration gates over 137 folds: all pass the location gate, the foundation and club-block samplers fail the lower-tail gate (28 and 22 to 23 of 137 folds below the scenario tenth percentile against a ceiling of 21), and a joint club-match sampler that draws appearance and residual from one historical club-week passes by one fold at the frozen seed only, with re-drawn seeds placing its expected count outside the band. We set out the data contract, the leakage controls, the component model, the single-week and multi-week integer programmes and the sampler, state the measured precision of the decision reading (90\% half-width 2.05 points, minimal detectable effect 3.49), and name every reading that was added after a result. The negative results are the finding: on this history the component line has reached the limit the decision reading can resolve, and the next gain has to come from information the archive does not hold or from the objective the solver reads, not from another estimator on the same features.
\end{abstract}

\section{Introduction}

A manager in Fantasy Premier League chooses, before each of 38 deadlines, a squad of fifteen Premier League players whose real-match events score points under published rules. The choice is constrained: a budget of 100.0 million at the start, at most three players from one club, exactly two goalkeepers, five defenders, five midfielders and three forwards, a legal formation for the eleven who start, a captain whose points count twice, and between deadlines a limited number of free transfers with a four-point charge for each extra one. Several chips, played once each in a half-season window, suspend one rule for one week. The score that matters is the realized total over the season, or the rank it buys in a league of friends.

The problem has two halves that the literature usually studies apart. One is forecasting: how many points will each player score next week? The other is decision: given a forecast, which legal squad, lineup and captain, and which transfers, maximise it? The two meet at a single interface, the column of expected points a solver reads. That interface makes a natural but rarely tested claim available for measurement: that a forecast which is more accurate at the player level produces squads that score more. This paper measures that claim on five seasons of public data, under one frozen protocol, with exact optimization under the official rules on both sides of every comparison.

The study grew out of an engineering programme rather than a hypothesis. The system we describe, SquadOpt, publishes weekly squad advice to the members of a private league and solves one-week plans in the member's browser. Its research arm measures each candidate model or policy on committed, walk-forward records before anything is served, and its rules for doing so, pre-registration, frozen folds, a measurements index in which every committed record has a row, were written to stop a forecasting improvement being promoted on a player-level metric alone. Over October 2026 that discipline produced five candidate records measured both ways, and the two ways disagreed. That disagreement is the subject.

\paragraph{Contributions:} We make four.
\begin{enumerate}
\item A component forecasting model for FPL, appearance $\times$ conditional points, with a leakage-safe feature contract on public data, a cold-start ladder and an inference-time availability rule, and a decision layer of two exact integer programmes (a single-week squad model and a multi-week transfer-and-chip planner) that are deterministic under a declared solver budget (Sections~\ref{sec:data} to~\ref{sec:opt}).
\item A paired measurement of five pre-declared estimator and feature candidates at two levels on the same folds: a points-error rule on 184 walk-forward folds, and realized squad points on 147 frozen decision folds under the official scoring rules. Player-level gains did not reach the squad, and the sign of the two readings disagreed for four of five candidates (Section~\ref{sec:results}).
\item A separation of three axes that one number usually conflates: point accuracy, distribution quality and decision value. A gated mixture improved the second decisively while moving neither of the others (Section~\ref{sec:dist}).
\item A calibration study of three scenario samplers for squad-level risk on two frozen gates, with a seed-sensitivity reading that turned a one-fold pass into a non-robust one, and a statement of the decision reading's measured precision (Sections~\ref{sec:calib} and~\ref{sec:precision}).
\end{enumerate}

We also report two negative planning-layer measurements on the same tables, a rank-conditional debiasing of expected points before the solve and a mean-reversion discount for the later weeks of a multi-week plan, because they test the same interface from the decision side.

\paragraph{What This Paper Is Not:} It is not a claim that the component model is the best FPL forecast available, and it is not a prospective live result: the live benchmark of the 2026-27 season starts at gameweek 7 under its own pre-registration and is read at the season's end. The repository behind this paper holds many earlier records from a first programme built on a deterministic form baseline: searches of that baseline's two parameters by grid and by Bayesian optimization, risk-aversion frontiers, planning-horizon and season-chain studies, transfer-discipline and chip-timing measurements, and the first calibration of the earlier scenario generator. They are not part of this study, for three reasons stated once here. Eleven of them were measured on a solver time budget that a later decision replaced with a deterministic one, so their numbers are not comparable with anything measured since; six multi-week-path records were drawn by a sampler whose club indexing has since been corrected and are not reproducible from the current code; and all of them measured a control that the component model has replaced. The one earlier record this paper does use is the frozen evaluation of the component model against the ridge control, because it is the comparator every later record is paired to.

\section{Related Work}
\label{sec:related}

Published work on fantasy football splits along the same seam as the problem. On the forecasting side, OpenFPL \cite{groos2025} combines position-specific random forests and gradient boosting on FPL and Understat histories and reports one-week RMSE of 1.517 against 2.136 for a last-five-games average on players returning 3 to 4 points, with a prospective window of seven gameweeks (2024-25 \gw{}32 to 38) and development folds split by club rather than by time. Its result is a player-level error, not a squad result, and its validation design (club folds, borrowed Understat values for 2020-21) does not transfer to a temporal test on the same archive. Gupta \cite{gupta2019} ensembles ARIMA and LSTM forecasts over three seasons to pick one fifteen by linear programming with no transfers, and reports only that the squad's forecast total of 3718 points was 87 too high, with no baseline. Bonello et al.\ \cite{bonello2019} blend statistics, betting and social signals and report about 11 points a week over their statistical comparator in 2018-19, with a fresh lineup each week and no real transfer penalty, so the figure is a system result under relaxed rules. Matthews, Ramchurn and Chalkiadakis \cite{matthews2012} built a Bayesian sequential manager that reached roughly the top percentile of 2.5 million managers in 2010-11 after training on 2009-10; it is the one published work that scores the whole sequential problem, and it predates the current rules by more than a decade.

The hierarchical Bayesian line on player and team performance is better developed for football than for the fantasy game. Egidi and Gabry \cite{egidi2018} model Italian fantasy ratings with player, team and position nesting and explicit missing appearances; Whitaker et al.\ \cite{whitaker2020} infer player event abilities from touch data and add them to the Baio and Blangiardo \cite{baio2010} team goal model, lifting an over/under AUC in every block while their predictive goal log-likelihood fell behind the team baseline; Egidi, Pauli and Torelli \cite{egidi2018b} combine historical team parameters with bookmakers' odds; Scholtes and Karakus \cite{scholtes2024} correct expected goals for player and position. Macri-Demartino, Egidi and Torelli \cite{macridemartino2025} pool adaptively over periods, which motivates change-aware shrinkage over a fixed half-life. None of these reports a decision gain under FPL rules, and most need event or tracking data the public archive does not carry. Rossi et al.\ \cite{rossi2018} forecast injuries from GPS training load with decision trees at about 50\% precision and 80\% recall; the inputs are private physical measurements, so the result bears on what a public-data availability model is missing rather than on what it can do.

On evaluation, Gneiting and Raftery \cite{gneiting2007} set out proper scoring rules and why coverage and sharpness alone do not determine forecast quality; our distribution reading (Section~\ref{sec:dist}) uses the logarithmic score and probability-integral-transform tails in that sense. Duan et al.\ \cite{duan2020} learn conditional distribution parameters by natural-gradient boosting on non-football benchmarks. For the inference on paired per-fold differences we use a moving-block bootstrap \cite{kunsch1989} confined to seasons, and for the precision reading the covariate-adjustment logic of Deng et al.\ \cite{deng2013}. For the solvers we use CP-SAT from OR-Tools \cite{ortools} on the server and HiGHS \cite{huangfu2018} in the browser; both are exact for the integer programmes of Section~\ref{sec:opt}. The estimators are scikit-learn's \cite{pedregosa2011}.

What the literature does not contain, as far as we found, is a paired reading of a forecasting change at both levels on the same folds, with the squad chosen by an exact solver under the official rules on both sides. That is the measurement this paper adds.

\section{The Game and the Decision Problem}
\label{sec:game}

A gameweek is a set of fixtures; most clubs play once, a club may play twice (a double gameweek) or not at all (a blank). A player's gameweek points are the sum over the club's fixtures of the official scoring rules: points for appearing (one for under sixty minutes, two for sixty or more), for goals and assists by position, for a clean sheet to a goalkeeper or defender who played sixty minutes, for saves and penalty saves, bonus points, and deductions for cards, own goals and conceded goals. From 2025-26 the rules add two points for a defensive-contribution threshold (ten clearances, blocks, interceptions and tackles for a defender, twelve such actions plus recoveries for a midfielder or forward) \cite{plglossary}. This change matters for the data below: 2025-26 points include a component that earlier seasons could not have.

A squad is fifteen players: two goalkeepers, five defenders, five midfielders, three forwards, at most three from one club, bought within the budget. Before each deadline the manager names eleven starters in a legal formation (one goalkeeper, three to five defenders, two to five midfielders, one to three forwards), orders the four on the bench and names a captain and vice-captain. The captain's points count twice; if the captain does not play, the vice-captain's do. If a starter plays no minutes, the first bench player who keeps the formation legal replaces him (automatic substitution). The manager may make transfers between deadlines: one free transfer accrues per week and up to five may be banked; each transfer beyond the free ones costs four points. Four chips, each once per half-season, lift a rule for one week: Bench Boost scores all fifteen, Triple Captain triples the captain, Wildcard allows unlimited free transfers, Free Hit allows them for one week and reverts the squad afterwards.

We write the single-week decision as a 0--1 programme in Section~\ref{sec:opt}. Its objective reads one number per player, the expected points $\hat y_i$ for the coming gameweek, so the forecasting model of Section~\ref{sec:model} enters the decision only through that column. Everything the model knows about variance, rotation, fixtures or dependence between players either reaches $\hat y_i$ or does not reach the squad at all. That is the structural fact the study exploits.

\section{Data}
\label{sec:data}

\subsection{The Archive and the Canonical Panel}

The historical data are the public FPL archive maintained by Vaastav Anand \cite{vaastav}, pinned at commit \texttt{8c97b2ad} and verified by checksum on every machine. The supported range is six seasons, 2020-21 to 2025-26, because the two earlier seasons omit position and club; it holds 156{,}075 player-gameweek rows for 1{,}960 players, 1{,}133 of whom appear in more than one season. Every row carries season, gameweek, player, club, position, price, minutes and total points, and the panel is keyed on (season, gameweek, player).

Four properties of the source shaped the adapter, and each was established by inspection rather than assumed. Player identity is the archive's \texttt{code}, which survives every season boundary (479 of 479 shared players), not its per-season \texttt{id} (1 of 479). Club identity is likewise the persistent club code, which survives all 85 consecutive-season club pairings where the per-season integer survives 53. Duplicate rows arise from two unlike causes and are treated unalike: a repeated fixture record is dropped, a genuine double gameweek is summed. The archive's price column is recorded at an undocumented time and differs from the roster's current cost for 537 of 692 players in 2025-26, so prices are shifted back one gameweek before use: a stale price costs accuracy, a leaky one costs correctness. Two small corrections are recorded: 101 goalkeeper rows were restored by accepting the spelling \texttt{GKP} beside \texttt{GK} in one gameweek of 2021-22, and the twenty manager entries per season introduced in 2024-25 are excluded as not squad-eligible.

Fixtures are stored at their own grain, one row per club per fixture, because a club with two fixtures in a gameweek has no single opponent or venue. Across the six seasons 189 club-gameweeks hold more than one fixture. The player-gameweek view derives from it by explicit aggregation to a fixture count and home and away counts; a fixture without a gameweek (a postponement awaiting a date) is excluded and counted, so a club genuinely has no fixture that week, and postponement history is unrecoverable rather than invented.

\subsection{Time of Knowledge and Leakage Control}

Every canonical column carries a time-of-knowledge class. Pre-match columns (price, club, position, fixture counts) may be read from the target row; outcome columns (minutes, points) may be read only from rows strictly before the target gameweek. All aggregation of outcome columns passes through one primitive that groups by (season, player), confirms ascending gameweek order within the group, and applies a shift of one before a rolling window. Because there is one such code path there is one place to audit, and the primitive refuses to roll a pre-match column (which would discard information) or an unclassified one. Rate features are ratios of shifted sums, not means of per-row ratios, and are left undefined, never zero, when the window holds no minutes. Insufficient history yields a missing value that is never back-filled.

Five test families enforce the structure: mutating outcomes in rows at or after $t$ leaves every feature of rows before $t$ bit-identical; features computed on a panel truncated at $t$ equal those from the full panel; a season's opening gameweeks are unaffected by the previous season's close; row order does not change output; and hand-computed rolling values are asserted literally. The same checks are made at fold level: a fold built from the full panel equals one built from a panel truncated after its decision gameweek, and rewriting the decision gameweek's outcome must move the realized score.

Three archive columns are deliberately excluded from every model. The archive's own expected-points column is scraped after the match. Ownership share has unverified snapshot timing. Availability status, chance of playing and news are recorded after the fact; their as-of time cannot be recovered, so a coefficient fitted on them would be fitted on information not held at the deadline. Availability enters only at live inference, as a multiplier after the projection (Section~\ref{sec:avail}). The source's fixture-difficulty rating is ingested for provenance and never consumed.

\subsection{Walk-Forward Folds}

A decision point is a season and gameweek. Rows strictly before it fit the model; rows through it, with the target gameweek's pre-match columns, build the projection; the target gameweek's realized points are read only after the decision is frozen. Random splits are not expressible: no function in the splitting layer accepts a seed, shuffle or fraction, and a test asserts the public surface never grows one. Each season's opening gameweek is skipped as a decision because a season-scoped rolling feature has no history there; the opener uses a separate carry-over and price-prior path that this paper does not measure.

Three fold populations appear below, nested by purpose (Table~\ref{tab:folds}). The \emph{player-level} population is 184 folds over 2021-22 to 2025-26, 130{,}095 scored rows, of which 128{,}627 take the component route and 1{,}468 (players without complete features, mostly first gameweeks) take the direct-control fallback. The \emph{decision} population is the 147 folds of 2021-22 to 2024-25 frozen when the component model was first evaluated against the ridge control, on its own 101{,}447-row table; it excludes 2025-26. The \emph{calibration} population is 137 of those 147: the first nine folds of 2021-22 are burn-in for the residual history the sampler needs, and one further fold is abstained because its frozen decision selected a direct-control player whose scenario has no component prediction. These populations were fixed before the candidates of Section~\ref{sec:candidates} existed; no fold was added or removed after a result.

\begin{table}[t]
\centering\small
\caption{The three fold populations. A fold is one decision gameweek. The player-level population reads 2025-26 as declared development data; the decision and calibration populations exclude it.}
\label{tab:folds}
\begin{tabularx}{\linewidth}{L R L L}
\toprule
Population & Folds & Seasons & Reads \\
\midrule
Player level (points-error rule) & 184 & 2021-22 to 2025-26 & 130{,}095 player-gameweek rows, 128{,}627 on the component route \\
Decision (realized squad points) & 147 & 2021-22 to 2024-25 & one solved squad per arm per fold, official scoring \\
Calibration (S1, S2 gates) & 137 & 2021-22 \gw{}11 to 2024-25 \gw{}38 & 1{,}000 scenarios per fold on the frozen decision \\
\bottomrule
\end{tabularx}
\end{table}

\subsection{The Status of 2025-26}

The repository locked 2025-26 as a final holdout in August 2026 and spent it once, under a pre-registered plan, on an earlier candidate that this paper does not report. Every record here that reads 2025-26 does so under a separately declared development scope, states so in its manifest (\texttt{locked\_holdout\_accessed: true}, \texttt{development\_scope: v2}), and is presented as development evidence, never as an unseen test. 2025-26 enters only the player-level rule. The decision and calibration readings end at 2024-25. The 2025-26 points target also includes the defensive-contribution points the earlier seasons lack (Section~\ref{sec:game}); nothing is added or re-thresholded to bridge the eras, and the manifests record the boundary.

\section{The Forecasting Model}
\label{sec:model}

\subsection{Component Decomposition}

For player $i$ at decision gameweek $t$ with $k_i$ scheduled fixtures, let $A_i \in \{0,1\}$ be whether he plays at least one minute, $M_i$ his minutes and $Y_i$ his points. The model estimates
\begin{equation}
p_i = \Pr(A_i = 1 \mid x_i),\qquad
\mu_i = \mathbb{E}[M_i \mid A_i = 1, x_i],\qquad
g_i = \mathbb{E}[Y_i \mid A_i = 1, x_i],
\end{equation}
and hands the solver
\begin{equation}
\hat y_i = p_i \, g_i, \qquad \hat m_i = p_i\,\mu_i,
\label{eq:compose}
\end{equation}
with the bounds $p_i \in [0,1]$, $\mu_i \in [0, 90k_i]$, $\hat y_i \ge 0$ applied after estimation. A blank gameweek ($k_i = 0$) sets all three to zero whatever the history says. A player who did not appear scores exactly zero, so non-appearance is an atom at zero rather than a small residual, and the decomposition keeps that atom separate from the conditional scoring rate. The pre-registration that fixed this structure also defines a start component ($\Pr(\text{started} \mid A=1)$) with start and substitute conditional means; the archive carries no verified start label across the supported seasons, so that component is recorded as unavailable and no proxy (sixty minutes, lineup membership) is substituted for it.

The decomposition exists because the two quantities fail differently. A first-choice striker who is injured and a fringe striker who is fit can carry the same recent points-per-ninety while having nothing in common as a selection; collapsing availability and rate into one regression makes the residual unattributable to either. It also matches the uncertainty layer (Section~\ref{sec:sampler}), where a player whose variance comes from rotation is a different object from one whose variance comes from finishing.

\subsection{Feature Contract}

The served contract (\texttt{v1}) has nine columns, all as-of the deadline: the shifted rolling means over the player's own season of minutes, points and appearance rate, each over windows of three and five gameweeks, plus the price in tenths, the club's fixture count in the target gameweek and its home-fixture count. Windows are over gameweeks, not fixtures, and use a minimum period of one, so a player's second gameweek already carries one prior observation; his first has none and takes the fallback route. The wider contract measured as a candidate (\texttt{v2}) adds the one- and ten-gameweek windows for minutes and appearance rate (thirteen columns); the points windows stay at three and five. Cross-season carry-over columns were tried in private research and left out of \texttt{v2} because they are missing for every player without a qualifying earlier season, which would send half the rows to the fallback under the no-imputation policy.

\subsection{Estimators}

The served model fits, inside each fold on rows strictly before its gameweek, three pipelines on standardized features: a logistic regression ($C = 1$, L-BFGS, 1{,}000 iterations) for appearance on all complete-feature rows, and two ridge regressions ($\alpha = 1$, Cholesky) for conditional minutes and conditional points on the rows that appeared. A fold needs at least 200 training rows for each fit; a thinner fold is refused and its rows take the fallback. The fallback (\emph{direct control}) is the frozen ridge control projection of the earlier programme, fitted walk-forward on the same panel with cross-season carry-over, so a row the component model cannot price still gets a leakage-safe number rather than zero. The scaler, the estimators and any weights are fitted inside the training slice; nothing is standardized on the whole panel.

The expected points the served model publishes are therefore a product of a probability and a conditional mean, both linear in standardized features through their link functions. The point of the candidates below is to ask whether more flexible estimators of the same three quantities, or more features, move the squads that product selects.

\subsection{Declared Candidates}
\label{sec:candidates}

Six candidates were declared under the same comparison rule (Section~\ref{sec:rule}), each in a method note or pre-registration committed before its measurement. The first, a season-age weighting of training rows by $0.5^{\text{age}}$, lost on the rule in September 2026 and its numbers were not committed; it is counted here as a comparison made but not reported. The five with committed records are:

\begin{description}
\item[Boosted pair.] Histogram gradient-boosted trees replace all three estimators: a classifier for appearance, squared-error regressors for minutes and points, each with learning rate 0.05, 300 iterations, 31 leaves, 100 minimum samples per leaf, L2 regularization 1.0, no early stopping, fixed seed. The constants were declared and not searched. The motivation came from private research on player-by-fixture rows (nine chronological windows, 28{,}117 shared rows) where trees on the same information beat the linear pair by 16\% on appearance log loss and minutes MAE.
\item[Wider feature contract.] The boosted pair on \texttt{v2} features, compared against the boosted pair on \texttt{v1}; the windows are the only difference.
\item[Points-only boosting.] The linear appearance and minutes estimators kept, the points regressor boosted. Declared after the boosted pair's decision reading suggested that its player-level gain lived in the appearance probability; its outcome on the decision folds was known from a recomposition probe before the declaration, and the record says so.
\item[Regime-gated experts.] One linear expert set per regime, the cross of fixture regime (single or double gameweek) and rule era (before or from the 2025-26 defensive-contribution rules), mixed by a multinomial logistic gate on the nine features plus three regime terms (double indicator, era indicator, season age). Each expert is the served pair on its regime's rows only; a regime with fewer than 200 rows is absent and the gate renormalizes.
\item[Gated mixture.] The appearance estimator kept; each conditional ridge replaced by a gated mixture of two ridge experts with a supervised regime label read from the realized training target: for minutes, a \emph{full} expert on rows with at least 60 minutes per fixture and a \emph{short} expert on the rest; for points, a \emph{return} expert on rows with at least four points and a \emph{no-return} expert on the rest. The gate is the served logistic estimator on the nine features; the expected value is the gate-weighted mean of the experts. This candidate also publishes a predictive distribution, the gate-weighted mixture of the experts' Gaussian residual distributions with each scale the expert's own residual standard deviation on its training rows (floored at one minute and half a point); the control's comparator distribution is the single Gaussian around the ridge mean.
\end{description}

\subsection{Availability as a Rule at Inference}
\label{sec:avail}

At a live deadline the system captures the roster, prices, fixtures and availability itself and stamps the capture instant, so availability becomes usable without being learned. A stated chance of playing sets a multiplier to that fraction; otherwise the status decides (available 1, injured, suspended, unavailable and not in squad 0, doubtful without a stated chance 0.5); a player with no record is treated as available; an unrecognised status stops the run. The multiplier is applied after the projection and reported beside it, players reduced to zero are reported rather than removed, and nothing in training reads historical availability. On the 2026-27 pre-season capture (584 players) the rule zeroed 56 and reduced 14. None of the measurements in this paper uses it: the archive cannot support it, and the walk-forward readings are therefore readings of the model without the live rule.

\section{Optimization}
\label{sec:opt}

\subsection{The Single-Week Squad Programme}

Let $I$ be the player pool, $I_p$ the players in position $p$, $I_c$ those of club $c$, $\pi_i$ the price in tenths and $\hat y_i$ the expected points. Binary variables $x_i$ (in the squad), $s_i$ (starts) and $c_i$ (captain); the bench is $x_i - s_i$. With squad quotas $Q_p$, formation bounds $[\ell_p, u_p]$, budget $B$, club cap 3 and bench weight $w_b$:
\begin{align}
\max\quad & \sum_i \hat y_i s_i + \sum_i \hat y_i c_i + w_b \sum_i \hat y_i (x_i - s_i) \label{eq:obj}\\
\text{s.t.}\quad & \sum_i x_i = 15,\qquad \sum_i s_i = 11,\qquad \sum_i c_i = 1, \nonumber\\
& \sum_{i \in I_p} x_i = Q_p,\qquad \ell_p \le \sum_{i \in I_p} s_i \le u_p \quad \forall p, \nonumber\\
& \sum_i \pi_i x_i \le B,\qquad \sum_{i\in I_c} x_i \le 3 \quad \forall c, \nonumber\\
& s_i \le x_i,\qquad c_i \le s_i \quad \forall i. \nonumber
\end{align}
The captain term adds his expectation a second time. The bench weight ($w_b = 0.1$ throughout) values squad depth without modelling automatic substitution; a starter receives no bench term. Bench order and the vice-captain are set after the solve by expected points and are not variables.

The programme is solved with CP-SAT \cite{ortools} on integers. Expected points are scaled by $10^3$ and rounded half-up through decimal arithmetic (the maximum rounding error is 0.0005 points a player; Python's built-in rounding is not used); prices are already integers. Players are ordered by identifier before variables are created, the solver runs one search worker under a fixed seed, and the measurement configuration stops the search on a \emph{deterministic} budget of 5.0 work units rather than on elapsed time, with a wall-clock cap that a run must not reach. An incumbent stopped by the deterministic budget is reproducible under the pinned solver version, which an elapsed-time limit cannot promise across machines. When the primary solve proves optimality, a second solve fixes the objective at its optimum and breaks ties by preferring lower stable ranks, captain first, then lineup, then squad; when it returns only a feasible incumbent, the tie-break is skipped. The result records the status, the best bound, the gap and whether the wall clock stopped the search. Extracted solutions are rechecked against every constraint before they are returned.

Solver status is part of the evidence. Under the 5.0 budget the linear component arm shared by all five candidate records ended 106 of 147 decision folds proven optimal and 41 feasible without proof; the ridge control 117 and 30. A feasible-but-unproven incumbent is reproducible but carries search noise in its realized score of unknown direction, because realized points are not the solver's objective. Section~\ref{sec:results} splits one comparison by this status.

\subsection{The Multi-Week Planner}

The planner optimizes a sequence of deadlines over a horizon $t = 1,\dots,H$ with a fixed player universe and supplied buy and sell prices (sell at or below buy; the optimizer does not infer the game's half-profit rule). Variables per player and week are $x_{it}, s_{it}, c_{it}$ as above plus $\mathrm{in}_{it}$ and $\mathrm{out}_{it}$; the squad, formation, club and captain constraints hold in every week. Continuity and bank accounting replace the single budget inequality:
\begin{equation}
x_{it} = x_{i,t-1} + \mathrm{in}_{it} - \mathrm{out}_{it},\quad \mathrm{in}_{it} + \mathrm{out}_{it} \le 1,\quad
b_t = b_{t-1} + \sum_i \sigma_{it}\,\mathrm{out}_{it} - \sum_i \beta_{it}\,\mathrm{in}_{it} \ge 0,
\end{equation}
with $x_{i,0}$ and $b_0$ the current squad and bank. Free transfers follow the game's accrual with exact max and min equalities: with $n_t = \sum_i \mathrm{in}_{it}$ and $f_t$ free transfers before deadline $t$, $\mathrm{paid}_t = \max(n_t - f_t, 0)$, $\mathrm{unused}_t = \max(f_t - n_t, 0)$ and $f_{t+1} = \min(5, \mathrm{unused}_t + 1)$. The objective discounts later weeks and charges hits at $h$ points each:
\begin{equation}
\max \sum_{t=1}^{H} \gamma^{\,t}\Big( \sum_i \hat y_{it} s_{it} + \sum_i \hat y_{it} c_{it} + w_b \sum_i \hat y_{it}(x_{it}-s_{it}) - h\,\mathrm{paid}_t \Big).
\end{equation}
No terminal bank value is added: money is useful only when it buys a better legal plan inside the horizon.

Chips are one Boolean per available chip per week, at most one chip a week and each at most once a horizon; the caller derives availability from the season's two half-season windows. Bench Boost adds the unweighted remainder $(1 - w_b)\hat y_{it}$ for every bench player through an auxiliary Boolean bounded by the chip, the squad variable and one minus the starter variable. Triple Captain adds the captain's expectation once more through an auxiliary bounded by the chip and the captain variable. Wildcard moves paid transfers to inequality form, $\mathrm{paid}_t \ge n_t - f_t - 15\,\mathrm{wc}_t$, so the charge vanishes under the chip, and the entering free-transfer total is preserved after it as the official rules state. Free Hit carries per-week base squad and bank variables pinned to this week's or the previous week's values by the chip's Boolean, so the following week starts from the squad the Free Hit week started from. A two-tier tie-break prefers fewer chips among equal plans, then the later week, since a rolling planner re-decides a deferred chip with fresher information. All discounted weights are scaled and rounded half-up, a weight that rounds to zero is refused, and extraction re-verifies the chip accounting and squad continuity. Projections for later weeks of the horizon repeat the decision week's projection scaled by that club's fixture count; a learned multi-week projection is not part of the served model.

\subsection{Where the Programmes Run}

The member website is static and solves one-week plans, chip values and rival-aware variants in the browser with HiGHS \cite{huangfu2018} compiled to WebAssembly, from a published device-plan table of about 600 players. Three- and five-week windows run on the server with CP-SAT. Two committed measurements compared the solvers on the live window model: on the 30 three- and five-week windows of fifteen members at gameweek 7 of 2026-27, CP-SAT proved 23 at the 5.0 budget and left 7 five-week windows feasible; native HiGHS proved all 30 within 1{,}800 s, three at CP-SAT's incumbent and four above it (by 0.06 to 21.08 points), all at or below CP-SAT's best bound, and six of the seven within CP-SAT's wall time. The WebAssembly build proved all 30 at 1{,}800 s but only 10 of 15 three-week and none of the five-week windows within 15 s, so windows stay on the server. Both the one-week programme on the device and the windows on the server are the same models as the measurement path; what differs is the solver and the time allowed.

\section{Scenario Sampler and Squad-Level Risk}
\label{sec:sampler}

The solver reads a mean. A member also wants to know how a squad's week can go, and the system's risk layer answers with scenarios: joint draws of every player's minutes and points for the coming gameweek, scored under the official rules into a distribution of squad scores. The sampler's inputs are the component model's out-of-fold table: for each history row, the predicted $p_i, \mu_i, g_i$ and the realized appearance, minutes and points.

\paragraph{Foundation Sampler:} For scenario $j$ and player $i$: draw $A_i^{(j)} \sim \text{Bernoulli}(p_i)$; if $A_i^{(j)} = 0$, minutes and points are exactly zero; if $A_i^{(j)} = 1$, take one historical row that appeared and add its paired residuals, $M_i^{(j)} = \text{clip}(\mu_i + \varepsilon^M, 1, 90k_i)$ and $Y_i^{(j)} = g_i + \varepsilon^Y$, where $(\varepsilon^M, \varepsilon^Y)$ are that row's realized minus predicted conditional minutes and points. Points are not clipped at zero, since an FPL score can be negative. One source fold (a past gameweek) is chosen per scenario, so every player's residual in a scenario comes from the same historical week; this is a block-bootstrap boundary and no more, since the residual rows within the fold are drawn independently per player. Residual history is strictly earlier than the target fold, needs at least eight such folds, and excludes rows without a component prediction. The sampler publishes the appearance matrix itself, so the scorer never infers appearance from minutes; the one-minute floor exists so that the published minutes cannot contradict the published appearance. Every draw is fingerprinted over all three matrices and its provenance.

\paragraph{Club-Block Sampler:} The same appearance draw, then for every club in the target squad one source club-week block inside the scenario's source fold, chosen among that fold's clubs with the same fixture count; each player's residual row is drawn within his club's block, among rows of his position where the block has any. A club that rotated in the source week lifts no co-movement into the appearance draw.

\paragraph{Joint Club-Match Sampler:} The appearance is drawn from the block too. With the source row's observed appearance $a_s$ and predicted probability $p_s$, its latent uniform is $u = v p_s$ if $a_s = 1$ and $u = p_s + v(1-p_s)$ if $a_s = 0$, for $v \sim U(0,1)$; the target appears iff $u < p_i$. Marginally $\Pr(A_i = 1) = p_i$ exactly when the source rows' appearance model is calibrated on the block; a target that appears from a source row that appeared takes that row's paired residual, otherwise a residual is re-drawn among the block's appeared rows. Zero atom and club co-movement thus come from one draw.

\paragraph{Calibration Gates:} For fold $t$ with frozen decision $d_t$, scenario scores $S_t^{(j)}$ under the official automatic-substitution and captain rules and realized score $R_t$:
\begin{equation}
\mathrm{PIT}_t = \tfrac{1}{1000}\sum_j \mathbf{1}[S_t^{(j)} \le R_t],\qquad q^{10}_t = \text{linear tenth percentile of } S_t,\qquad L_t = \mathbf{1}[R_t < q^{10}_t].
\end{equation}
Gate S1 (location): the mean PIT over the 137 folds lies in $[0.43, 0.57]$. Gate S2 (lower tail): the mean of $L_t$ lies in $[0.04, 0.16]$, which on 137 folds is at most 21 folds. Both gates were frozen before the first sampler was measured, with 1{,}000 scenarios a fold and seed 0; no shift, scale or recentering is applied, and the point estimates decide. The population and both gates are shared by every sampler record below.

\section{Evaluation Protocol}
\label{sec:protocol}

\subsection{The Player-Level Rule}
\label{sec:rule}

Every candidate is exported as an out-of-fold table from the same commit as its control, with the same rows, targets, routes and folds. The comparison script refuses two arms that are not one paired measurement. Per fold, the metric is the mean absolute error of the composed expected points $\hat y_i$ against realized gameweek points over all scored rows; the paired statistic is control MAE minus candidate MAE, so a positive value favours the candidate. The rule, declared on 7 September 2026 and unchanged since: the candidate is preferred only if (i) its 2025-26 points MAE is lower, (ii) the 90\% season-aware moving-block bootstrap interval of the paired per-fold differences on 2025-26 lies entirely above zero, and (iii) the same interval on 2024-25 does not lie entirely below zero. No minimum effect size is imposed. Appearance Brier and log loss, conditional minutes MAE and row coverage are reported beside the rule and do not enter it. The primary and control-era seasons are constants of the script, not options.

\subsection{The Decision-Level Reading}

For each of the 147 decision folds, each arm's expected points go through the programme of Section~\ref{sec:opt} under the measurement configuration (one worker, seed 0, deterministic budget 5.0), the squad is completed (bench order and vice-captain by expected points), and the squad is scored against the fold's realized outcomes under the official automatic-substitution and captain rules (the vice-captain replaces a non-playing captain; a non-playing starter is replaced by the first bench player who keeps the formation legal). Roster, prices, positions, outcomes, solver limits and scorer are identical between arms; only the expected-points column differs. The reading reports each arm's mean realized score, the paired candidate-minus-control mean and median, wins, ties and losses, and a 90\% interval. Zero-minute starters and automatic-substitution points are reported as diagnostics.

\subsection{Inference}

Intervals on paired per-fold differences come from a moving-block bootstrap \cite{kunsch1989} that resamples blocks of four consecutive folds within a season and never across a season boundary, 5{,}000 resamples, seed 0, percentile 90\% interval. The repository's threshold for \emph{certifying} a points effect is a mean of at least 1.5 points a gameweek with an interval excluding zero and at most one losing season; that is a declared bar, not a measured noise level, and the three pin conditions declared with the candidates here ask only for the rule to pass and a non-negative decision mean.

\subsection{Pre-Registration, and Where It Fell Short}

Each candidate has a declaration committed before its measurement: method notes for the boosted pair, the wider contract and points-only boosting, pre-registrations for the regime and mixture candidates and for both sampler candidates. The method notes precede the measurement, not the code; the regime, mixture and sampler protocols precede the code. The decision-level reading was declared with the rule for points-only boosting, the regime experts and the mixture; for the boosted pair and the wider contract it was added after the first record and is reported beside the rule, flagged as such in the records and in Table~\ref{tab:decision}. The points-only decision reading reproduces a recomposition probe known before its declaration, so it is not an independent confirmation. The joint sampler's seed reading was added after its result. The regime gate read twelve inputs where its declaration said eleven; this was recorded and the declaration left as written. We report these as they fell because a protocol section that omitted them would be the one claim in the paper the records contradict.

\subsection{Multiple Comparisons}

Six candidates were measured under one rule, and no correction is applied. The study is exploratory; its confirmation is the prospective live benchmark of 2026-27, read from gameweek 7 under its own pre-registration. For the decision reading the question is moot in one direction: no interval excluded zero before any correction.

\section{Results}
\label{sec:results}

\subsection{Player Level}

Table~\ref{tab:player} gives the pooled player-level metrics over the 184 folds and the rule's two seasons. Two candidates were preferred. The boosted pair lowered 2025-26 points MAE from 1.0058 to 0.9872, with the paired interval \ci{+0.0143}{+0.0212} above zero and 2024-25 \ci{+0.0068}{+0.0169} also above; it was better pooled on all three components (appearance log loss 0.3462 to 0.3308, minutes MAE 16.16 to 15.41). The wider contract, compared against the boosted pair, lowered 2025-26 points MAE again to 0.9679 (\ci{+0.0150}{+0.0240}) and the other components with it (0.3154, 14.39). The other three were not selected: points-only boosting left appearance and minutes unchanged by construction and was worse on 2025-26 points MAE with an interval below zero (\ci{-0.0041}{-0.0005}); the regime experts were worse on appearance and points and lower only on minutes; the mixture was equal on appearance, better on minutes in every season and worse on points in both rule seasons.

\begin{table}[t]
\centering\small
\caption{Player-level metrics on the 184 walk-forward folds (pooled over 128{,}627 component-route rows) and the pre-declared rule on 2025-26 and 2024-25 points MAE. Differences are reference minus candidate per fold with 90\% season-aware block-bootstrap intervals; positive favours the candidate. The wider contract's reference is the boosted pair; every other reference is the linear control.}
\label{tab:player}
\setlength{\tabcolsep}{3pt}
\footnotesize
\resizebox{\linewidth}{!}{%
\begin{tabular}{l rrr r r r l}
\toprule
& \multicolumn{3}{c}{Pooled} & \multicolumn{2}{c}{2025-26 points MAE} & 2024-25 & \\
\cmidrule(lr){2-4}\cmidrule(lr){5-6}\cmidrule(lr){7-7}
Arm & App.\ LL & Min.\ MAE & Pts MAE & MAE & diff.\ (90\%) & diff.\ (90\%) & Verdict \\
\midrule
Linear control & 0.3462 & 16.161 & 1.0716 & 1.0058 & & & reference \\
Boosted pair & 0.3308 & 15.406 & 1.0550 & 0.9872 & $+0.0185$ \ci{+0.0143}{+0.0212} & $+0.0118$ \ci{+0.0068}{+0.0169} & preferred \\
Wider contract$^{b}$ & 0.3154 & 14.386 & 1.0400 & 0.9679 & $+0.0194$ \ci{+0.0150}{+0.0240} & $+0.0103$ \ci{+0.0079}{+0.0127} & preferred \\
Points-only boosting & 0.3462 & 16.161 & 1.0714 & 1.0079 & $-0.0021$ \ci{-0.0041}{-0.0005} & $-0.0030$ \ci{-0.0055}{+0.0014} & no selection \\
Regime-gated experts & 0.3464 & 16.104 & 1.0729 & 1.0179 & $-0.0122$ \ci{-0.0162}{-0.0080} & $-0.0017$ \ci{-0.0028}{-0.0001} & no selection \\
Gated mixture & 0.3462 & 16.038 & 1.0733 & 1.0061 & $-0.0004$ \ci{-0.0012}{+0.0004} & $-0.0020$ \ci{-0.0029}{-0.0009} & no selection \\
\bottomrule
\end{tabular}}
\par\smallskip\footnotesize $^{b}$Compared against the boosted pair, not the linear control.
\end{table}

\subsection{Decision Level}

Table~\ref{tab:decision} and Figures~\ref{fig:folds} and~\ref{fig:seasons} give the decision reading. The linear components score 63.42 realized points a gameweek against the ridge control's 58.52, a paired $+4.89$ (the frozen first evaluation, under a ten-second time limit on an earlier export, read $+4.62$ with 84 wins, 8 ties and 55 losses and a positive mean in every season). Against that linear arm, no candidate separates. The two candidates the rule preferred score less on the mean: the boosted pair $-0.75$ (\ci{-2.66}{+1.05}, 75/5/67) and the wider contract $-0.95$ (\ci{-3.24}{+1.01}, 67/8/72; against the boosted pair $-0.20$, \ci{-1.80}{+1.06}). Two of the three the rule rejected score more: points-only boosting $+0.47$ (\ci{-1.19}{+2.22}, 73/6/68) and the regime experts $+0.54$ (\ci{-0.63}{+1.67}, 39/86/22). The mixture scores $-0.36$ (\ci{-1.07}{+0.73}, 33/78/36). All five paired means lie inside 1.5 points in size; four of the five intervals reach past 1.5 on one side. The regime and mixture candidates, which share the linear appearance estimator, tie the control on more than half the folds (86 and 78 of 147): their squads are often identical.

\begin{table}[t]
\centering\small
\caption{Realized squad points a gameweek on the 147 frozen decision folds, each arm's squads solved by the same programme and scored under the official rules. Paired differences are candidate minus linear components with 90\% season-aware block-bootstrap intervals and wins/ties/losses. ``After'' marks a decision reading added after the player-level result rather than declared with the rule.}
\label{tab:decision}
\setlength{\tabcolsep}{5pt}
\begin{tabular}{l rr r l r l}
\toprule
Arm & Mean & vs ridge & vs linear & 90\% interval & W/T/L & Declared \\
\midrule
Ridge control & 58.52 & & & & & \\
Linear components & 63.42 & $+4.89$ & & & & reference \\
Boosted pair & 62.67 & $+4.14$ & $-0.75$ & \ci{-2.66}{+1.05} & 75/5/67 & after \\
Wider contract & 62.46 & $+3.94$ & $-0.95$ & \ci{-3.24}{+1.01} & 67/8/72 & after \\
Points-only boosting & 63.88 & $+5.36$ & $+0.47$ & \ci{-1.19}{+2.22} & 73/6/68 & with rule$^{a}$ \\
Regime-gated experts & 63.96 & $+5.44$ & $+0.54$ & \ci{-0.63}{+1.67} & 39/86/22 & with rule \\
Gated mixture & 63.05 & $+4.53$ & $-0.36$ & \ci{-1.07}{+0.73} & 33/78/36 & pre-registered \\
\midrule
\multicolumn{7}{l}{\footnotesize Frozen first evaluation, ten-second limit: ridge 58.54, linear components 63.16, paired $+4.62$, 84/8/55.} \\
\bottomrule
\end{tabular}
\par\smallskip\footnotesize $^{a}$Declared with the rule, but its value reproduces a recomposition probe known before the declaration.
\end{table}

\begin{figure}[t]
\centering
\includegraphics[width=\linewidth]{fig1_decision_fold_differences}
\caption{Per-fold paired decision differences, candidate minus linear components, by season, for the five candidates on the 147 decision folds. Each point is one gameweek; the horizontal mark is the season mean. The spread within a season is of the order of ten points either way; the means are within one.}
\label{fig:folds}
\end{figure}

\begin{figure}[t]
\centering
\includegraphics[width=0.85\linewidth]{fig2_decision_by_season}
\caption{Season means of the paired decision difference for each candidate. No candidate is above the control in every season; the sign changes across seasons for all five.}
\label{fig:seasons}
\end{figure}

\paragraph{Solver Status:} The points-only record splits its reading by the linear arm's solver status: candidate minus linear is $+0.78$ on the 106 folds the linear arm proved optimal and $-0.34$ on the 41 it left feasible. The split does not change the verdict, but it shows that part of a paired difference of this size can be search noise on unproven incumbents, which is why the status is recorded on every solve.

\paragraph{Zero-Minute Starters:} The boosted pair's squads started more players who then played no minutes (108 against the linear arm's 83 over 147 folds), despite a lower appearance log loss. Its appearance probabilities were sharper, and the composition $p_i g_i$ hands the solver a product in which a confident probability on a high-rate player outweighs a diffuse one on a steady starter. This is the one diagnostic that points at a mechanism; the records state it as a hypothesis, and no record measures it.

\subsection{Accuracy Against Decision Value}

Figure~\ref{fig:scatter} plots each candidate's 2025-26 points-MAE difference against its decision difference. Four of five lie off the diagonal: the two preferred by the rule are below the control at the decision level, two rejected by the rule are above it. The mixture alone has the same sign on both, and it is worse on both. On this history, the points-error rule and realized squad points are not monotone in each other for the changes measured.

\begin{figure}[t]
\centering
\includegraphics[width=0.62\linewidth]{fig3_points_mae_against_squad_points}
\caption{Points-MAE improvement on the rule season (horizontal; positive favours the candidate) against the paired decision difference on the 147 folds (vertical). The rule's verdict and the decision sign disagree for four of five candidates.}
\label{fig:scatter}
\end{figure}

\subsection{Distribution Quality}
\label{sec:dist}

The mixture publishes a predictive distribution per player; the control's comparator is a single Gaussian around the ridge mean. Both were scored as discretised distributions on the 52{,}517 rows that appeared and took the component route in both arms, with the logarithmic score $\log\big(F(y+\tfrac12)-F(y-\tfrac12)\big)$ and the mid-point PIT. Table~\ref{tab:dist} gives the pooled reading. For points the mean log score rose from $-2.495$ to $-2.043$, the mean PIT moved from 0.465 to 0.505, and both tails moved toward their nominal 0.05 (lower 0.004 to 0.027, upper 0.074 to 0.041). For minutes the log score rose from $-4.725$ to $-4.182$, the lower tail moved from 0.068 to 0.045, and the upper tail moved from 0.033 to 0.016, further from nominal. On the 74 folds scored fold by fold (2024-25 and 2025-26) the mixture won every fold for both components; the per-fold mean differences were $+0.38$ (\ci{+0.36}{+0.40}) and $+0.51$ (\ci{+0.49}{+0.52}) for points and $+0.52$ and $+0.50$ for minutes. The same candidate was no better as a mean (Table~\ref{tab:player}) and no better at the decision (Table~\ref{tab:decision}). Distribution quality moved without the other two axes moving.

\begin{table}[t]
\centering\small
\caption{Predictive-distribution scores of the gated mixture and its single-Gaussian control, pooled over 52{,}517 appeared rows. Higher log score is better; tails are nominally 0.05.}
\label{tab:dist}
\begin{tabular}{llrrrr}
\toprule
Component & Arm & Mean log score & Mean PIT & Lower tail & Upper tail \\
\midrule
Points & control & $-2.495$ & 0.465 & 0.004 & 0.074 \\
Points & mixture & $-2.043$ & 0.505 & 0.027 & 0.041 \\
Minutes & control & $-4.725$ & 0.506 & 0.068 & 0.033 \\
Minutes & mixture & $-4.182$ & 0.510 & 0.045 & 0.016 \\
\bottomrule
\end{tabular}
\end{table}

\subsection{Scenario Calibration}
\label{sec:calib}

Table~\ref{tab:calib} and Figure~\ref{fig:pit} give the three samplers on the two gates. Every arm passes S1: mean PIT lies between 0.515 and 0.521 in every record, so the scenario distributions are located where the realized scores fall. The lower tail is where they differ. The foundation sampler puts the realized score below its own scenario tenth percentile in 28 of 137 folds (0.204) in the binding record and in the joint record, and 25 in the block record's regenerated table; the ceiling is 21. The club-block sampler reads 22 in its own record (0.161, one fold outside) and 23 in the joint record. The joint club-match sampler reads 21 (0.153) at the frozen seed 0, inside the band by one fold; its paired lower-tail flag against the foundation is $-0.051$ (\ci{-0.066}{-0.022}) and against the block sampler $-0.015$ (\ci{-0.037}{0.000}).

\begin{table}[t]
\centering\small
\caption{Calibration of the three samplers on the frozen 137 folds, 1{,}000 scenarios a fold, seed 0. S1 band $[0.43, 0.57]$ on the mean PIT; S2 band $[0.04, 0.16]$ on the share of folds with the realized score below the scenario tenth percentile, at most 21 folds.}
\label{tab:calib}
\begin{tabular}{llrrrl}
\toprule
Sampler & Record & Mean PIT & Below $q^{10}$ & Rate & S2 \\
\midrule
Foundation & binding (Phase D) & 0.5147 & 28 & 0.204 & fail \\
Foundation & block record, control arm & 0.5197 & 25 & 0.182 & fail \\
Club block & block record, candidate & 0.5166 & 22 & 0.161 & fail \\
Foundation & joint record, control A & 0.5209 & 28 & 0.204 & fail \\
Club block & joint record, control B & 0.5186 & 23 & 0.168 & fail \\
Joint club-match & joint record, candidate & 0.5207 & 21 & 0.153 & pass (seed 0) \\
\bottomrule
\end{tabular}
\end{table}

\begin{figure}[t]
\centering
\includegraphics[width=\linewidth]{fig4_sampler_pit}
\caption{PIT histograms of the three samplers over the 137 calibration folds. All three are centred; the excess mass at the lowest bin is the lower-tail failure that S2 measures.}
\label{fig:pit}
\end{figure}

The pass is not robust. A reading added after the result re-drew the twelve folds whose PIT lies near 0.10 under seeds 0 to 9 (Figure~\ref{fig:seeds}). The joint sampler's count ranged from 21 to 25 with mean 23.4, inside the band on one seed of ten; the block sampler's ranged from 23 to 26 with mean 24.0, inside on none. The plug-in expected count for the joint sampler is 22.7 to 22.9, above the ceiling. The within-club appearance co-movement the sampler was built to carry is almost absent in the measured squads (correlation 0.0030 against $-0.0007$ for the block sampler), and its lower tenth percentile against the block sampler is mostly a location shift of about $-0.16$ points rather than a fatter tail. The 137 folds had been read by three samplers before this one, and the joint protocol was committed after the block record had listed the folds it flipped. The record therefore stands as development evidence on reused folds. Nothing was pinned: the live risk layer still carries no calibrated sampler version, and the repository's rule for one requires a reading on folds no sampler was chosen on.

\begin{figure}[t]
\centering
\includegraphics[width=\linewidth]{fig5_s2_counts}
\caption{Lower-tail fold counts of the joint and club-block samplers over seeds 0 to 9 against the S2 ceiling of 21 folds. The frozen seed 0 is the only seed on which the joint sampler is inside the band.}
\label{fig:seeds}
\end{figure}

What the three readings agree on is informative in itself. Location is right in every sampler, and every sampler is too narrow in the lower tail by roughly a factor of 1.3 to 1.9 in the count of tail events. Club co-movement, the mechanism the two candidates added, moved the count from 28 toward 21 to 25; what remains is of the same order. The projection's own uncertainty, which none of the samplers carries (each draws residuals around the point forecast as if it were known), is the plausible remainder.

\subsection{Planning-Layer Measurements}
\label{sec:planning}

Two measurements on the same tables test the interface from the decision side (Table~\ref{tab:planning}). The first asks whether the winner's curse, the upward bias of the selected players' projections, can be corrected before the solve: a rank-conditional additive debias fitted walk-forward on the boosted table's own residuals, applied to expected points before the programme, against the same table unshrunk, on 176 folds over five seasons with every solve proven. The debiased arm scored $-0.85$ points a gameweek (\ci{-1.59}{-0.02}, 51/63/62), below the control in four of five seasons; a position-level shrinkage control arm read $+0.07$ (\ci{-0.28}{+0.56}). The diagnostic explains the result: the control's own selection gap on this table is $+0.51$ a selected starter, where the earlier form baseline had measured $+2.95$, so there was little curse left to correct and the bucket shifts moved the decision only at bucket boundaries.

The second asks whether the churn of a rolling three-week planner can be damped by shrinking the later weeks' projections toward the position mean while keeping their level, $\hat y_{i,k} = \max\big(0,\, m_{\text{pos}} + \rho^{k}(\hat y_i - m_{\text{pos}})\big)\cdot k_i$ for horizon offset $k$. Against the naive three-week planner on 184 weeks, the discounted planner read $-0.56$ (\ci{-1.51}{+0.89}) at $\rho = 0.75$, $-0.03$ at $0.9$ and $-0.14$ at $0.6$; every three-week arm read below the one-week control, the naive one by $-0.51$ (\ci{-1.79}{+0.34}). Neither measurement certified anything, and both say the same thing as the forecasting candidates from the other side: on this history, changes to what the solver reads move the realized score by less than the reading can resolve.

\begin{table}[t]
\centering\small
\caption{Two planning-layer measurements, realized points a gameweek net of charged hits, arm minus comparator, 90\% season-aware block-bootstrap intervals. Both ``measured, not certified'' under their pre-registered gates.}
\label{tab:planning}
\begin{tabular}{llrrl}
\toprule
Measurement & Arm vs comparator & Folds & Mean & 90\% interval \\
\midrule
Rank-conditional shrinkage & rank debias vs unshrunk & 176 & $-0.85$ & \ci{-1.59}{-0.02} \\
 & position shrinkage vs unshrunk & 176 & $+0.07$ & \ci{-0.28}{+0.56} \\
Mean-reversion discount & 3-week $\rho{=}0.75$ vs 3-week naive & 184 & $-0.56$ & \ci{-1.51}{+0.89} \\
 & 3-week $\rho{=}0.9$ vs 3-week naive & 184 & $-0.03$ & \ci{-0.83}{+0.64} \\
 & 3-week $\rho{=}0.6$ vs 3-week naive & 184 & $-0.14$ & \ci{-0.76}{+0.90} \\
 & 3-week naive vs 1-week & 184 & $-0.51$ & \ci{-1.79}{+0.34} \\
 & 3-week $\rho{=}0.75$ vs 1-week & 184 & $-1.07$ & \ci{-2.52}{+0.26} \\
\bottomrule
\end{tabular}
\end{table}

\subsection{What the Decision Reading Can Resolve}
\label{sec:precision}

The precision of the decision reading was measured once, on the 147 paired component-minus-ridge differences, as a descriptive instrument study. Treating the folds as independent, the 90\% half-width of the paired mean is 2.05 points a gameweek and the minimal detectable effect at $\alpha = 0.05$ and power 0.8 is 3.49. The strongest pre-decision covariate, the total projected points of the component pool, correlates with the paired difference at $r = 0.38$ and would reduce those to 1.89 and 3.23 if an external expectation for it were justified; with the covariate centred on the same folds the adjusted mean equals the raw mean, so no precision gain is claimed. Folds overlap in squads and weeks within a season, so the independence figures are a lower bound on the true width. Every candidate effect in Table~\ref{tab:decision} is below one half-width in size. The study was, in that sense, powered to find effects of three points a gameweek and found none of one.

\section{Discussion}
\label{sec:discussion}

\paragraph{Why Accuracy Does Not Reach the Squad:} Three features of the interface explain the disagreement without appeal to chance. First, the solver reads a mean, and the mean enters it as the composed product $p_i g_i$. A sharper appearance probability lowers log loss and points MAE because most realized points are zero and a confident zero is cheap; the same sharpness hands the solver a product that can rank a volatile high-rate player above a steady starter. The boosted pair's extra zero-minute starters are that effect's footprint. Second, the decision is a selection of fifteen from about six hundred under tight constraints, and it is decided by the ordering near the top of each position's price band, not by the error over all rows; two forecasts can differ in MAE by a few hundredths and agree on the fifteen players the constraints admit. The regime and mixture candidates produced identical squads to the control on more than half the folds. Third, realized squad points are noisy at roughly 16 points a gameweek fold to fold, so a paired change of under one point needs several hundred folds to resolve at the usual levels; five seasons give 147.

\paragraph{Three Axes, Not One:} The mixture result separates what a single ``accuracy'' usually bundles. A forecast can be sharper and better calibrated as a distribution without moving its mean, and a mean can move without moving the decision. For the risk layer the first axis is what matters; for the squad the third; the points-error rule sits on the second and predicts neither. A system that promotes on one metric will be surprised by the others, and the surprises here ran in both directions.

\paragraph{Where the Next Gain Is:} The negative results constrain the search. On the same nine as-of features, four estimator families and one wider contract moved the decision by less than the reading resolves, and the two planning-side corrections did likewise. The repository's own composition probes, not reported as records, tried the boosted points component with linear appearance ($+0.47$ against linear) and the reverse ($-0.79$), and an appearance exponent that lowered points MAE and moved decisions nowhere. What has not been tried, because the archive does not hold it, is information: pre-deadline availability and projected line-ups with verifiable as-of times, and market odds for the match-level quantities (clean sheets, goals) that drive the position-level returns. The live capture now records availability at the deadline, so the prospective season can measure the first of these. The other open direction is the objective. The solver maximises a sum of means; a member in a league of fifteen is maximising rank, which is a different functional of the same scenarios, and the calibration work of Section~\ref{sec:calib} is the precondition for reading it.

\paragraph{On the Calibration Result:} We report the joint sampler's seed reading as the finding rather than its seed-0 pass. A one-fold margin with 1{,}000 draws a fold is inside Monte Carlo noise (the standard error of an empirical tenth percentile at that sample size is near 0.01), and a gate read on reused folds after seeing which folds flip is development evidence by construction. The method that stops that error is cheap: re-draw at several seeds and state the expected count before calling a gate passed, and list every candidate already read on the same folds. We did the second after the first result rather than before it, and record that.

\section{Limitations}
\label{sec:limitations}

\begin{itemize}
\item \textbf{2025-26 is development data.} It enters the player-level rule after having been read by earlier work; no reading here is an unseen test. Its points include the defensive-contribution component the four earlier seasons lack, so the rule season and the decision seasons are scored under different rules.
\item \textbf{The decision folds end at 2024-25} and were fixed when the component model was first evaluated. They were not chosen for the candidates, but neither were they drawn afresh, and the 137 calibration folds have now been read by four samplers.
\item \textbf{Not every solve is proven.} Under the 5.0 deterministic budget, 41 of the linear arm's 147 decision solves and 30 of the ridge control's ended feasible without proof. Their incumbents are reproducible, but their realized scores carry search noise of unknown direction. The points-only split by status is the only quantification.
\item \textbf{Readings added after results.} Two decision readings, one seed reading, one probe reproduction and three bootstrap identifiers were added after a result; one declaration precedes its measurement but not its outcome. Section~\ref{sec:protocol} names each.
\item \textbf{Six comparisons, no correction.} The player-level verdicts are exploratory; the live 2026-27 benchmark is the confirmation.
\item \textbf{No odds, no line-ups, no event data.} Every feature derives from the public player-gameweek archive and the fixture list. The related work that reports strong results uses richer inputs.
\item \textbf{One archive, one game.} The archive's price timing is undocumented, postponements are unrecoverable, and the start label is absent, so the start component is unmeasured. Nothing here is claimed for other fantasy games.
\item \textbf{The availability rule is unmeasured.} It runs live and cannot be validated on the archive; the walk-forward readings are of the model without it.
\end{itemize}

\section{Reproducibility}
\label{sec:repro}

Every number in Sections~\ref{sec:results} is extracted by a script from eleven committed JSON records, each with its commit, manifest, input digests and declaration, and nothing in the extraction fits or measures. From a clean clone of the repository, \texttt{python -m scripts.paper\_record\_tables} writes the tables and prints ``83 of 83 record checks hold''; \texttt{python -m scripts.paper\_figures} draws the five figures from those tables; a check mode compares every output's SHA-256 against the package manifest. The records themselves were produced by the repository's runners from the archive at commit \texttt{8c97b2ad}, which is not redistributed but is pinned and checksummed; reproducing a record means re-running its runner under its pre-registration, which the records document command by command. Eleven of the twelve regenerated outputs match their recorded digests byte for byte; the two that name the checked-out commit differ in that field only. The solvers are OR-Tools CP-SAT (one worker, seed 0, deterministic budget 5.0) and HiGHS 1.15.1; the estimators are scikit-learn 1.9 under Python 3.11 and 3.13 with pinned constraints. Related-work quotations were checked against the source documents line by line.

\paragraph{Use of Language Models:} The repository's maintainers use language-model assistants for code, for extracting tables from records and for review. This manuscript was drafted by one such assistant from the committed records, the pre-registrations and the method notes at the corresponding author's request; every number was taken from the records through the extractor, and the authors are responsible for the argument, the interpretation of every number and the text as submitted.

\section{Conclusion}

We measured five pre-declared changes to a component forecasting model for Fantasy Premier League at two levels on the same walk-forward folds: a points-error rule, and realized squad points when the squad is chosen by an exact integer programme under the official rules. Two changes were better forecasts by the rule; neither produced better squads, and the sign of the two readings disagreed for four of the five. A mixture model improved the predictive distribution on every fold scored while moving neither the mean nor the decision. Three scenario samplers were located correctly and all too narrow in the lower tail; the one that passed the tail gate did so by one fold at one seed. Two planning-side corrections on the same tables also failed to separate from their controls. The decision reading is precise to about two points a gameweek on 147 folds, and no measured effect reached one.

The engineering conclusion is the paper's: a forecast change for a constrained selection problem has to be read at the decision, with the same solver on both sides and the solver's status recorded, before it is promoted on a player-level metric; and the three axes of a forecast, point accuracy, distribution quality and decision value, have to be read separately because they moved separately here. The scientific one is narrower and, we think, useful: on five seasons of public FPL data and nine as-of features, the component line has reached what the decision reading can resolve, and the next measurable gain has to come from information the archive does not hold or from the objective the solver reads.

\paragraph{Acknowledgements:} The public archive is maintained by Vaastav Anand. The league members whose weekly decisions the system serves are the reason the measurements exist.

\begin{thebibliography}{99}\small

\bibitem{groos2025} D.\ Groos. OpenFPL: An open-source forecasting method rivaling state-of-the-art Fantasy Premier League services. Author preprint, arXiv:2508.09992, 2025.

\bibitem{gupta2019} A.\ Gupta. Time series modeling for Dream Team in Fantasy Premier League. arXiv:1909.12938, 2019.

\bibitem{bonello2019} N.\ Bonello, J.\ Beel, S.\ Lawless and J.\ Debattista. Multi-stream data analytics for enhanced performance prediction in fantasy football. In \emph{Proceedings of AICS}, 2019. arXiv:1912.07441.

\bibitem{matthews2012} T.\ Matthews, S.\ Ramchurn and G.\ Chalkiadakis. Competing with humans at fantasy football: team formation in large partially-observable domains. In \emph{Proceedings of the 26th AAAI Conference on Artificial Intelligence}, pages 1394--1400, 2012.

\bibitem{egidi2018} L.\ Egidi and J.\ Gabry. Bayesian hierarchical models for predicting individual performance in football (soccer). \emph{Journal of Quantitative Analysis in Sports}, 2018. DOI 10.1515/jqas-2017-0066.

\bibitem{whitaker2020} G.\ A.\ Whitaker, R.\ Silva and D.\ Edwards. A Bayesian inference approach for determining player abilities in football. arXiv:1710.00001, revised 2020.

\bibitem{baio2010} G.\ Baio and M.\ Blangiardo. Bayesian hierarchical model for the prediction of football results. \emph{Journal of Applied Statistics}, 37(2):253--264, 2010.

\bibitem{egidi2018b} L.\ Egidi, F.\ Pauli and N.\ Torelli. Combining historical data and bookmakers' odds in modelling football scores. \emph{Statistical Modelling}, 18(5--6):436--459, 2018.

\bibitem{scholtes2024} A.\ Scholtes and O.\ Karakus. Bayes-xG: player and position correction on expected goals using Bayesian hierarchical approach. arXiv:2311.13707, 2023; published 2024.

\bibitem{macridemartino2025} R.\ Macri-Demartino, L.\ Egidi and N.\ Torelli. Adaptive period pooling for football score models. arXiv:2508.05891, 2025; \emph{Journal of the Royal Statistical Society, Series C}, 2026.

\bibitem{rossi2018} A.\ Rossi, L.\ Pappalardo, P.\ Cintia, F.\ M.\ Iaia, J.\ Fernández and D.\ Medina. Effective injury forecasting in soccer with GPS training data and machine learning. \emph{PLOS ONE}, 13(7):e0201264, 2018.

\bibitem{dixon1998} M.\ J.\ Dixon and M.\ E.\ Robinson. A birth process model for association football matches. \emph{Journal of the Royal Statistical Society, Series D}, 47(3):523--538, 1998.

\bibitem{gneiting2007} T.\ Gneiting and A.\ E.\ Raftery. Strictly proper scoring rules, prediction, and estimation. \emph{Journal of the American Statistical Association}, 102(477):359--378, 2007.

\bibitem{duan2020} T.\ Duan, A.\ Avati, D.\ Y.\ Ding, K.\ K.\ Thai, S.\ Basu, A.\ Y.\ Ng and A.\ Schuler. NGBoost: natural gradient boosting for probabilistic prediction. In \emph{Proceedings of ICML}, 2020.

\bibitem{kunsch1989} H.\ R.\ Künsch. The jackknife and the bootstrap for general stationary observations. \emph{Annals of Statistics}, 17(3):1217--1241, 1989.

\bibitem{deng2013} A.\ Deng, Y.\ Xu, R.\ Kohavi and T.\ Walker. Improving the sensitivity of online controlled experiments by utilizing pre-experiment data. In \emph{Proceedings of WSDM}, 2013.

\bibitem{ortools} L.\ Perron and F.\ Didier. CP-SAT, OR-Tools. Google, 2025. \url{https://developers.google.com/optimization/cp/cp_solver}.

\bibitem{huangfu2018} Q.\ Huangfu and J.\ A.\ J.\ Hall. Parallelizing the dual revised simplex method. \emph{Mathematical Programming Computation}, 10(1):119--142, 2018. HiGHS, \url{https://highs.dev}.

\bibitem{pedregosa2011} F.\ Pedregosa et al. Scikit-learn: machine learning in Python. \emph{Journal of Machine Learning Research}, 12:2825--2830, 2011.

\bibitem{plglossary} Premier League. FPL glossary: fantasy talk explained. 2026. \url{https://www.premierleague.com/en/news/2683145}.

\bibitem{vaastav} V.\ Anand. Fantasy Premier League historical data archive. \url{https://github.com/vaastav/Fantasy-Premier-League}, commit \texttt{8c97b2adb123863c3dd581e730f1360e89815ac2}.

\end{thebibliography}

\appendix

\section{The Committed Records}

Table~\ref{tab:records} lists the records the results read, with the pull request that merged each and the declaration it was measured under. All are under \texttt{docs/} in the repository.

\begin{table}[h]
\centering\small
\caption{Records behind the tables and figures. R1 to R11 as the reproducibility package names them.}
\label{tab:records}
\begin{tabularx}{\linewidth}{l L l}
\toprule
Id & Record & Declaration \\
\midrule
R1 & Phase C component evaluation (frozen, 147 folds, ridge control) & component model pre-registration \\
R2 & Phase D component squad calibration (binding, foundation sampler) & Phase D calibration pre-registration \\
R3 & Mean-reversion discount & its pre-registration \\
R4 & Gated mixture & its pre-registration \\
R5 & Regime-gated experts & its pre-registration \\
R6 & Joint club-match sampler & its pre-registration \\
R7 & Wider feature contract & method note \\
R8 & Boosted pair & method note \\
R9 & Points-only boosting & method note \\
R10 & Rank-conditional shrinkage & its pre-registration \\
R11 & Club-block sampler & its pre-registration \\
\bottomrule
\end{tabularx}
\end{table}

\section{Notation}

\begin{tabularx}{\linewidth}{l L}
\toprule
Symbol & Meaning \\
\midrule
$A_i, M_i, Y_i$ & appearance indicator, minutes and points of player $i$ in the target gameweek \\
$p_i, \mu_i, g_i$ & appearance probability, conditional minutes, conditional points \\
$\hat y_i = p_i g_i$ & expected points handed to the solver \\
$k_i$ & the club's fixture count in the target gameweek \\
$x_i, s_i, c_i$ & in squad, starts, captain (binary) \\
$w_b$ & bench weight, 0.1 \\
$\mathrm{PIT}_t, q^{10}_t, L_t$ & probability integral transform, scenario tenth percentile and lower-tail flag of fold $t$ \\
S1, S2 & the location and lower-tail calibration gates \\
\bottomrule
\end{tabularx}

\end{document}
